\Subhead{4.3.2}{Шумы лазеров. Формула Таунса и предельная стабильность частоты}

Спонтанное излучение добавляет к полю случайные комплексные приращения. Стабилизация амплитуды подавляет радиальные возмущения, но автономная фаза не имеет восстанавливающей силы и диффундирует. Если $\langle[\phi(t+\tau)-\phi(t)]^2\rangle=2D_\phi|\tau|$, гауссовское усреднение даёт
\begin{equation}
 g^{(1)}(\tau)=\e^{-D_\phi|\tau|},\qquad
 \Delta\omega_{\rm FWHM}=2D_\phi,\quad \Delta\nu_{\rm FWHM}=D_\phi/\pi.
\end{equation}
Это лоренцева линия. Численный множитель требует явных определений ширины, времени жизни и модели шума. В простой одномодовой модели с числом фотонов $N$, скоростью спонтанного добавления фотонов в моду $R_{\rm sp}$ и изотропными квадратурными приращениями получаем
\begin{equation}\label{eq:schawlow}
 2D_\phi=\frac{R_{\rm sp}}{2N},\qquad
 \Delta\nu=\frac{R_{\rm sp}}{4\pi N}
 \simeq\frac{n_{\rm sp}\kappa}{4\pi N}.
\end{equation}
Множитель можно восстановить из нормировки комплексной амплитуды $\alpha=\sqrt N\e^{\ii\phi}$. При $\langle|\dd\alpha|^2\rangle=R_{\rm sp}\dd t$ изотропный шум делится поровну между квадратурами, а $\dd\phi\simeq\operatorname{Im}(\e^{-\ii\phi}\dd\alpha)/\sqrt N$. Поэтому $\langle(\dd\phi)^2\rangle=R_{\rm sp}\dd t/(2N)=2D_\phi\dd t$.

В последнем равенстве~\eqref{eq:schawlow} $\kappa=1/\tau_p$ — скорость потери фотонов, $R_{\rm sp}=n_{\rm sp}\kappa$ в принятой эффективной модели. При одинаковых коэффициентах вынужденных переходов и населённостях $N_2>N_1$ фактор $n_{\rm sp}=N_2/(N_2-N_1)\geq1$. Предполагаются $N\gg1$ и стабилизированная амплитуда. Коэффициент подтверждается формулой FWHM $\Delta\omega=n_{\rm sp}/(2N\tau_p)$ у Rana~\cite{rana}, PDF-слайд~37. Так как $P_{\rm out}=\hbar\omega\kappa_{\rm out}N$, ширина убывает обратно мощности и зависит от полных и выходных потерь. Речь о пределе выбранной модели, а не о гарантированной линии реального лазера.

Амплитудно-фазовая связь в полупроводнике даёт дополнительный фактор типа $1+\alpha_H^2$; технические шумы накачки, длины и температуры могут доминировать. Ширина стационарной линии и стабильность частоты за заданное время усреднения — разные показатели. Дрейф резонатора не устраняется одним увеличением мощности; нужен внешний эталон или стабилизация.

% Источниковое расхождение: Kärtner, полный PDF-с.327 / печатная с.319,
% даёт CW-формулу Delta f=Theta/(2 pi N tau_p). Множитель два не согласован
% прочитанным фрагментом; этот источник не используется как подтверждение
% численного коэффициента eq:schawlow. Rana, слайд 37, совпадает с нашим выводом.
